\documentclass[]{article}

\usepackage{amsmath}
\usepackage{multirow}
\usepackage{natbib}
\usepackage{graphicx} 


\begin{document}



In this study, we investigated the statistical nature of Lagrangian vortex transport in three-dimensional isothermal turbulence to determine if the motion of these coherent structures follows classical diffusion or exhibits anomalous behavior. To address this, we analyzed a comprehensive dataset of vortex trajectories extracted from a high-resolution direct numerical simulation. We employed a robust pipeline to identify vortices using the Q-criterion and track their vorticity-weighted centroids over time, enabling a direct statistical characterization of their motion.

Our analysis, based on the Mean Squared Displacement (MSD), establishes that the transport is unequivocally superdiffusive, with a diffusive exponent of $\alpha \approx 1.82$. This value, significantly greater than one, indicates a transport mechanism far more efficient than a simple random walk. To uncover the underlying process, we analyzed the Velocity Autocorrelation Function (VACF), which revealed a slow decay, signifying strong temporal correlations and a persistent 'memory' in the vortex velocity. Concurrently, the probability distribution of trajectory step sizes was found to be nearly Gaussian, allowing us to rule out Lévy flights as the driving mechanism. These two findings collectively support the conclusion that vortex superdiffusion is best described as a correlated random walk.

We further explored the physical origins of this behavior and found the transport to be anisotropic, with different diffusive exponents for each Cartesian direction. This anisotropy was directly linked to a similarly anisotropic coupling with the background fluid velocity; the direction with the strongest superdiffusion corresponded to the direction of the strongest correlation between vortex and fluid velocities. However, the motion is not entirely slaved to the background flow. An analysis of the motion relative to the vortex's own orientation, defined by the local vorticity vector, showed that vortices preferentially drift in the plane perpendicular to their axis. This suggests that the observed transport is a complex interplay between advection by the large-scale flow and the intrinsic dynamics of the vortex structures themselves. The robustness of these findings was confirmed through sensitivity tests on the vortex identification threshold and trajectory length.

In summary, this work provides a quantitative characterization of Lagrangian vortex transport in 3D turbulence. We have learned that these fundamental structures do not behave as passive tracers undergoing classical diffusion. Instead, their motion is a robustly superdiffusive, correlated random walk. This behavior arises from the persistence of their velocity, driven by a combination of strong advection by the surrounding flow and their own intrinsic dynamics. These findings contribute to a more complete physical picture of turbulent transport, highlighting that the coherent, non-local nature of vortices plays a crucial role in the efficient mixing of momentum and energy in turbulent flows.

\end{document}
                